\documentclass[12pt,a4paper]{article}
\usepackage{amsmath,amssymb}
\usepackage{graphicx}
\usepackage{hyperref}
\usepackage{geometry}
\usepackage{booktabs}
\usepackage{natbib}
\geometry{margin=2.5cm}
\hypersetup{colorlinks=true,linkcolor=blue,citecolor=blue,urlcolor=blue}

\begin{document}

\title{RIFT Cannot Replace Dark Matter:\\
Exhaustive Test of Kinetic Modifications to the Locked Action}

\author{Christopher Robert Barrick Wilson\\
\small Independent Researcher\\
\small \texttt{C.Isomorph@gmail.com}}
\date{\today}
\maketitle

\begin{abstract}
A recurring question for any modified-gravity theory is whether the
scalar degree of freedom can source galactic dark matter halos and
reproduce flat rotation curves without invoking collisionless dark matter
particles.  We subject the locked RIFT action
($\Lambda_0=0.003$, $m_0=0.01$~Mpc$^{-1}$) to three independent tests
against the NGC~3198 rotation curve ($v_{\rm flat}=150$~km/s):
(i)~static Yukawa profile (Sim~99): the $\Psi$ halo falls exponentially
    steeper than the $\rho\propto r^{-2}$ isothermal required for flat
    curves, with a density deficit of $2.7\times10^6\times$;
(ii)~time-domain PDE condensation (Sim~100): the galactic curvature
    coupling $2\Lambda_0|R_{\rm gal}|\approx5\times10^{-9}$~kpc$^{-2}$
    is $5\times10^7\times$ weaker than $m_0^2$; curvature condensation at
    galactic densities is impossible; attractive self-coupling
    ($\lambda<0$) produces tachyonic blow-up, not a stable halo;
(iii)~Vainshtein screening (Sim~103): the nonlinear kinetic term
    $(\alpha_V/4M_V^2)(\partial\Psi)^4$ remains deep in the linear
    regime for all scanned $(\alpha_V, M_V)$; the Vainshtein ratio
    $\mathcal{V}=\alpha_V(\delta\Psi')^2/M_V^2\leq3.7\times10^{-6}$.
All three mechanisms fail for the same root reason: the galactic field
gradient $|\delta\Psi'|\sim6\times10^{-8}$~kpc$^{-1}$, set by the weak
$\Lambda_0 R_{\rm gal}$ source, is six orders of magnitude below any
nonlinear threshold.  A direct matter coupling $\beta\sim1600\,\Lambda_0$
would close the density gap but is not contained in the locked action and
would violate cosmological homogeneity constraints.  The conclusion is
definitive: RIFT at its cosmological best-fit cannot replace dark matter
through any modification of the existing Lagrangian.  New first-principles
physics coupling $\Psi$ to baryonic matter is the minimum requirement.
\end{abstract}

\section{Introduction}

Flat rotation curves of spiral galaxies --- the canonical evidence for
dark matter --- require a halo density $\rho\propto r^{-2}$ (isothermal)
within the optical disk.  For NGC~3198, with $v_{\rm flat}=150$~km/s
and $R_d=2.72$~kpc, the required halo density at $r=10$~kpc is
\begin{equation}
\rho_{\rm iso}(10~{\rm kpc}) = \frac{v_{\rm flat}^2}{4\pi G r^2}
  = 4.2\times10^3\,M_\odot\,{\rm kpc}^{-3}\,.
\label{eq:iso}
\end{equation}

The RIFT scalar field $\Psi$ couples to spacetime curvature through the
non-minimal term $\Lambda_0\Psi R$ in the action.  Could this coupling
produce a self-gravitating $\Psi$ condensate that mimics a dark matter halo?
We test this exhaustively.

\section{Test I: Static Yukawa Profile (Sim~99)}

The static field equation for $\Psi$ in spherical symmetry (with the
FLRW background subtracted) gives a Yukawa profile:
\begin{equation}
\delta\Psi(r) = A\,\frac{e^{-m_0 r}}{r}
  + B\,\frac{e^{+m_0 r}}{r}\,.
\end{equation}
The physical (decaying) solution with $m_0^{-1}=10{,}000$~kpc gives an
effective energy density
$\rho_\Psi \propto (\nabla\delta\Psi)^2 \propto e^{-2m_0 r}/r^4$.
This falls exponentially steeply --- orders of magnitude faster than
$r^{-2}$ --- and cannot reproduce flat rotation curves for any value of
the amplitude $A$.

\paragraph{G$_{\rm eff}$ scan.} A scan over $G_{\rm eff}/G\in[1, 10^4]$
confirms: boosting $\Lambda_0$ changes the amplitude but not the profile
shape, which remains Yukawa.  No static solution produces
$\rho\propto r^{-2}$.

\section{Test II: Time-Domain Condensation (Sim~100)}

We integrate the $\Psi$ Klein--Gordon PDE on a fixed galactic background
over $10^4$~Mpc/$c$ ($\sim 10^{10}$~yr), testing whether:
(a) the galactic curvature source $2\Lambda_0|R_{\rm gal}|$ drives a
    condensate, or
(b) an attractive self-coupling $\lambda\Psi^3/3!$ stabilises a halo.

\paragraph{Curvature coupling.}
The galactic Ricci scalar $|R_{\rm gal}|\approx 8\pi G\rho_{\rm disk}
\approx 10^{-9}$~kpc$^{-2}$.  The coupling strength
$2\Lambda_0|R_{\rm gal}|\approx5\times10^{-9}$~kpc$^{-2}$ compared
to $m_0^2=10^{-4}$~kpc$^{-2}$: the source is $5\times10^7\times$
weaker than the restoring force.  No condensate forms.

\paragraph{Attractive self-coupling.}
A scan over $\lambda<0$ shows the field diverges tachyonically before
forming a stable halo at any galactic density.  The static (Sim~99) and
time-domain (Sim~100) limits agree, confirming adiabaticity and ruling out
history-dependence as a resolution.

\section{Test III: Vainshtein Screening (Sim~103)}

The Vainshtein mechanism introduces the nonlinear kinetic term
$(\alpha_V/4M_V^2)(\partial\Psi)^4$, enhancing the effective energy
density:
\begin{equation}
\rho_\Psi(r) = \frac{c^2}{8\pi G}\cdot\tfrac{1}{2}
  \bigl(\delta\Psi'\bigr)^2
  \left[1 + \underbrace{\frac{\alpha_V}{M_V^2}(\delta\Psi')^2}_{\mathcal{V}}\right].
\end{equation}
In the Poisson-limit ($m_0 r\ll 1$) the field gradient is:
\begin{equation}
\frac{d\delta\Psi}{dr} = \frac{2G\Lambda_0 M_{\rm enc}(r)}{r^2}
  \approx 6.1\times10^{-8}\,{\rm kpc}^{-1}\quad(r=10~{\rm kpc})\,.
\end{equation}

\begin{table}[h]
\centering
\caption{Vainshtein ratio at $r=10$~kpc for NGC~3198.}
\begin{tabular}{ccc}
\toprule
$\alpha_V$ & $M_V$ [kpc$^{-1}$] & $\mathcal{V}=\alpha_V(\delta\Psi')^2/M_V^2$ \\
\midrule
$1$    & $0.001$ & $3.7\times10^{-9}$ \\
$100$  & $0.001$ & $3.7\times10^{-7}$ \\
$1000$ & $0.001$ & $3.7\times10^{-6}$ \\
$1000$ & $0.01$  & $3.7\times10^{-8}$ \\
\bottomrule
\end{tabular}
\end{table}

Every combination is deep in the linear regime ($\mathcal{V}\ll1$).
The scalar energy density is
$\rho_\Psi(10~{\rm kpc}) = 1.6\times10^{-3}\,M_\odot\,{\rm kpc}^{-3}$,
a factor of $2.7\times10^6$ below the isothermal requirement
(eq.~\ref{eq:iso}).

\section{Root Cause and Required New Physics}

All three tests fail for the same reason.
The field gradient $|\delta\Psi'|\sim6\times10^{-8}$~kpc$^{-1}$ is
determined by the coupling strength $\Lambda_0 R_{\rm gal}$, which
is fixed by the cosmological best-fit value $\Lambda_0=0.003$ and the
galactic curvature $R_{\rm gal}\sim10^{-9}$~kpc$^{-2}$.
This is an unavoidable consequence of the locked action: the same
$\Lambda_0$ that ensures cosmological consistency makes the galactic
source negligible.

To close the $2.7\times10^6\times$ density gap, a direct matter coupling
$\beta\Psi\rho_b$ with $\beta\approx4.9\approx1600\,\Lambda_0$ would be
needed.  Such a coupling:
\begin{enumerate}
\item Is not derivable from the locked action;
\item Introduces a large new free parameter;
\item Would drive $\mathcal{O}(1)$ perturbations in the cosmological
      background, violating CMB homogeneity.
\end{enumerate}

\section{Conclusion}

RIFT at its cosmological best-fit cannot replace dark matter through any
modification of the existing Lagrangian.
The scalar field couples too weakly to galactic curvature to source
the required $\rho\propto r^{-2}$ halo, and no kinetic enhancement
(Yukawa profile, temporal condensation, Vainshtein nonlinearity) changes
this conclusion.
A first-principles coupling between $\Psi$ and baryonic matter ---
derivable from a symmetry principle rather than added by hand --- is
the minimum new physics required for RIFT to address the dark matter problem.

\end{document}
